% Generated by task tex-libraries from dossiers/lib-arklib.md, section A (and F.2, F.3 for the
% new pin). Snippets are copied from the dossier, which read them with `git show` at dca90385.
\begingroup\sloppy
\section{ArkLib: the objects, from scratch}\label{sec:gen-lib-ark-objects}

ArkLib is the library of interactive oracle reductions in which the blueprint states the proof
system: a protocol is a schedule of prover messages and verifier challenges, a reduction is a
prover and a verifier for it, and the security notions are predicates on verifiers. This section
introduces its objects in dependency order. \textbf{Every snippet is from ArkLib at the old pin
\code{dca90385}}, read with \code{git show} at that commit.

\paragraph{At the new pin \texorpdfstring{\code{7653a901}}{7653a901}.} The objects below exist with
the same shapes (checked by name): \lean{OracleReduction}, \lean{OracleVerifier.toVerifier},
\lean{KnowledgeStateFunction}, \lean{rbrKnowledgeSoundnessWorstCaseWith}, \lean{GuardedForm},
\lean{OracleReduction.append}, the guarded completeness theorems and
\lean{Verifier.}\allowbreak\lean{GuardedForm.}\allowbreak\lean{ofEmpty}. Every security definition is restated in VCVio's new
probability notation (\cref{sec:gen-lib-vcvio-new}): the \code{diff} of
\file{Security/}\allowbreak\file{RoundByRound.lean} between the pins is the notation and the lemma names
(\lean{probEvent_*} $\to$ \lean{prEvent_*}, \lean{OptionT.prEvent_mk_*}); the mathematical content
of \lean{KnowledgeStateFunction}, \lean{rbrKnowledgeSoundnessWorstCase(With)} and the two
implications between them is unchanged. ArkLib now calls this whole framework \enquote{legacy}
(\cref{sec:gen-lib-ark-legacy}).

\paragraph{Reading conventions.} A \lean{Type} is a set of values; \lean{Fin n} is
$\{0, \dots, n-1\}$; a \lean{structure} is a record; a \lean{class} is a record found
automatically by type; \lean{Prop} is a proposition; \lean{OracleComp spec α} (VCVio) is a
computation that may ask questions to the oracles named by \lean{spec} and returns an \lean{α};
\lean{ProbComp} is \lean{OracleComp} over one uniform-sampling oracle; \lean{OptionT m α} adds a
failure outcome to \lean{m α}; \lean{Pr[p | c]} is the probability that the outcome of \lean{c}
satisfies \lean{p} (failure is not an outcome; VCVio at \code{f9dc47d9}); \lean{StateT σ m}
threads a state of type \lean{σ}. \Cref{sec:gen-lib-vcvio} introduces VCVio's objects.

\subsection{Protocols and transcripts}\label{sec:gen-lib-ark-protocols}

\begin{define}{\lean{Direction}: who speaks in a round}
A round is either a prover message (\lean{P_to_V}) or a verifier challenge (\lean{V_to_P}).
\end{define}
\begin{leancode}
inductive Direction where
  | P_to_V  -- Message
  | V_to_P -- Challenge
deriving DecidableEq, Inhabited, Repr
\end{leancode}
\src{\file{ArkLib/}\allowbreak\file{OracleReduction/}\allowbreak\file{Prelude.lean:66-69}, ArkLib \code{dca90385}}

\begin{define}{\lean{ProtocolSpec n}: the schedule of an $n$-round protocol}
For each of the $n$ rounds, who speaks and the type of what is sent.
\end{define}
\begin{leancode}
@[ext]
structure ProtocolSpec (n : ℕ) where
  /-- The direction of each message in the protocol. -/
  dir : Fin n → Direction
  /-- The type of each message in the protocol. -/
  «Type» : Fin n → Type
deriving Inhabited
\end{leancode}
\src{\file{ArkLib/}\allowbreak\file{OracleReduction/}\allowbreak\file{ProtocolSpec/}\allowbreak\file{Basic.lean:30-36}, ArkLib \code{dca90385}}

\begin{define}{Message and challenge indices, messages, challenges}
\lean{MessageIdx} and \lean{ChallengeIdx} are the rounds of each direction, as subsets of
\lean{Fin n}; \lean{Message} and \lean{Challenge} are their types.
\end{define}
\begin{leancode}
/-- Subtype of `Fin n` for the indices corresponding to messages in a protocol specification -/
@[reducible, simp]
def MessageIdx (pSpec : ProtocolSpec n) :=
  {i : Fin n // pSpec.dir i = Direction.P_to_V}

/-- Subtype of `Fin n` for the indices corresponding to challenges in a protocol specification -/
@[reducible, simp]
def ChallengeIdx (pSpec : ProtocolSpec n) :=
  {i : Fin n // pSpec.dir i = Direction.V_to_P}
\end{leancode}
\src{\file{ArkLib/}\allowbreak\file{OracleReduction/}\allowbreak\file{ProtocolSpec/}\allowbreak\file{Basic.lean:50-58}, ArkLib \code{dca90385}}

\begin{leancode}
@[reducible, inline, specialize, simp]
def Message (pSpec : ProtocolSpec n) (i : MessageIdx pSpec) := pSpec.«Type» i.val

/-- Unbundled version of `Message`, which supplies the proof separately from the index. -/
@[reducible, inline, specialize, simp]
def Message' (pSpec : ProtocolSpec n) (i : Fin n) (_ : pSpec.dir i = .P_to_V) := pSpec.«Type» i

/-- The type of the `i`-th challenge in a protocol specification -/
@[reducible, inline, specialize, simp]
def Challenge (pSpec : ProtocolSpec n) (i : ChallengeIdx pSpec) := pSpec.«Type» i.val
\end{leancode}
\src{\file{ArkLib/}\allowbreak\file{OracleReduction/}\allowbreak\file{ProtocolSpec/}\allowbreak\file{Basic.lean:68-77}, ArkLib \code{dca90385}}

\begin{define}{\lean{FullTranscript}, \lean{Transcript k}: all entries, and a prefix}
\lean{FullTranscript} is the list of all $n$ entries; \lean{Transcript k} is its prefix of length
$k$ (the full transcript is \lean{Transcript (Fin.last n)} definitionally).
\end{define}
\begin{leancode}
@[reducible, inline, specialize]
def FullTranscript (pSpec : ProtocolSpec n) := (i : Fin n) → pSpec.«Type» i
\end{leancode}
\src{\file{ArkLib/}\allowbreak\file{OracleReduction/}\allowbreak\file{ProtocolSpec/}\allowbreak\file{Basic.lean:104-105}, ArkLib \code{dca90385}}

\begin{leancode}
@[reducible, inline, specialize]
def Transcript (k : Fin (n + 1)) (pSpec : ProtocolSpec n) : Type :=
  (pSpec⟦:k.val⟧).FullTranscript
\end{leancode}
\src{\file{ArkLib/}\allowbreak\file{OracleReduction/}\allowbreak\file{ProtocolSpec/}\allowbreak\file{Basic.lean:261-263}, ArkLib \code{dca90385}}

\begin{define}{Concatenation of schedules}
\lean{pSpec₁ ++ₚ pSpec₂} runs the first schedule then the second; \lean{seqCompose} iterates it
over a family.
\end{define}
\begin{leancode}
/-- Appending two `ProtocolSpec`s -/
abbrev append (pSpec : ProtocolSpec m) (pSpec' : ProtocolSpec n) : ProtocolSpec (m + n) :=
  ⟨Fin.vappend pSpec.dir pSpec'.dir, Fin.vappend pSpec.Type pSpec'.Type⟩

@[inherit_doc]
infixl : 65 " ++ₚ " => ProtocolSpec.append
\end{leancode}
\src{\file{ArkLib/}\allowbreak\file{OracleReduction/}\allowbreak\file{ProtocolSpec/}\allowbreak\file{SeqCompose.lean:38-43}, ArkLib \code{dca90385}}

\begin{leancode}
@[inline]
def seqCompose {m : ℕ} {n : Fin m → ℕ} (pSpec : ∀ i, ProtocolSpec (n i)) :
    ProtocolSpec (Fin.vsum n) where
  dir := Fin.vflatten (fun i => (pSpec i).dir)
  «Type» := Fin.vflatten (fun i => (pSpec i).Type)
\end{leancode}
\src{\file{ArkLib/}\allowbreak\file{OracleReduction/}\allowbreak\file{ProtocolSpec/}\allowbreak\file{SeqCompose.lean:358-362}, ArkLib \code{dca90385}}

\subsection{Oracles}\label{sec:gen-lib-ark-oracleinterface}

\begin{define}{\lean{OracleInterface M}: how a value of type \lean{M} is queried}
A query type and an answer, given as an oracle context over \lean{ReaderM M} (the value is read,
the answer computed). \lean{Response q} and \lean{answer m q} are derived. leanerVM's instance for
a column has queries in $\E^n$ and answers the multilinear extension
(\cref{sec:gen-l0-oracle}).
\end{define}
\begin{leancode}
@[ext]
class OracleInterface (Message : Type u) where
  Query : Type v
  toOC : OracleContext Query (ReaderM Message)
\end{leancode}
\src{\file{ArkLib/}\allowbreak\file{OracleReduction/}\allowbreak\file{OracleInterface.lean:54-57}, ArkLib \code{dca90385}}

\begin{leancode}
def Response {Message : Type*} [O : OracleInterface Message]
    (q : O.Query) : Type _ :=
  O.toOC.spec q

def spec {Message : Type*} [O : OracleInterface Message] :
    OracleSpec O.Query :=
  O.toOC.spec

@[implicit_reducible]
def answer {Message : Type*} [O : OracleInterface Message]
    (m : Message) (q : O.Query) : O.Response q :=
  (O.toOC.impl q).run m
\end{leancode}
\src{\file{ArkLib/}\allowbreak\file{OracleReduction/}\allowbreak\file{OracleInterface.lean:67-78}, ArkLib \code{dca90385}}

\subsection{Provers, verifiers, reductions}\label{sec:gen-lib-ark-reductions}

\begin{define}{\lean{Prover}: the honest prover as a state machine}
A state type per round; \lean{input} from statement and witness; \lean{sendMessage} and
\lean{receiveChallenge} per round (computations over the shared oracle \lean{oSpec});
\lean{output}, the final statement and witness. Provers of the soundness games are the same
structure (with any \lean{WitIn}, \lean{WitOut}).
\end{define}
\begin{leancode}
@[ext]
structure ProverState (n : ℕ) where
  PrvState : Fin (n + 1) → Type
\end{leancode}
\src{\file{ArkLib/}\allowbreak\file{OracleReduction/}\allowbreak\file{Basic.lean:132-134}, ArkLib \code{dca90385}}

\begin{leancode}
@[ext]
structure ProverRound {ι : Type} (oSpec : OracleSpec ι) {n : ℕ} (pSpec : ProtocolSpec n)
    extends ProverState n where
  /-- Send a message and update the prover's state -/
  sendMessage (i : MessageIdx pSpec) :
    PrvState i.1.castSucc → OracleComp oSpec (pSpec.Message i × PrvState i.1.succ)
  /-- Receive a challenge and update the prover's state -/
  receiveChallenge (i : ChallengeIdx pSpec) :
    PrvState i.1.castSucc → OracleComp oSpec (pSpec.Challenge i → PrvState i.1.succ)
\end{leancode}
\src{\file{ArkLib/}\allowbreak\file{OracleReduction/}\allowbreak\file{Basic.lean:154-162}, ArkLib \code{dca90385}}

\begin{leancode}
@[ext]
structure Prover {ι : Type} (oSpec : OracleSpec ι)
    (StmtIn WitIn StmtOut WitOut : Type)
    {n : ℕ} (pSpec : ProtocolSpec n) extends
      ProverState n,
      ProverInput StmtIn WitIn (PrvState 0),
      ProverRound oSpec pSpec,
      ProverOutput oSpec (StmtOut × WitOut) (PrvState (Fin.last n))
\end{leancode}
\src{\file{ArkLib/}\allowbreak\file{OracleReduction/}\allowbreak\file{Basic.lean:222-229}, ArkLib \code{dca90385}}

\begin{define}{\lean{Verifier}: a verdict on a full transcript}
A function of the statement and the full transcript, returning the output statement or failing
(rejecting).
\end{define}
\begin{leancode}
@[ext]
structure Verifier {ι : Type} (oSpec : OracleSpec ι)
    (StmtIn StmtOut : Type) {n : ℕ} (pSpec : ProtocolSpec n) where
  verify : StmtIn → FullTranscript pSpec → OptionT (OracleComp oSpec) StmtOut
\end{leancode}
\src{\file{ArkLib/}\allowbreak\file{OracleReduction/}\allowbreak\file{Basic.lean:247-250}, ArkLib \code{dca90385}}

\begin{define}{\lean{OracleProver}}
A prover whose input statement also carries the contents of the input oracles and whose output
carries the output oracles.
\end{define}
\begin{leancode}
@[reducible, inline]
def OracleProver {ι : Type} (oSpec : OracleSpec ι)
    (StmtIn : Type) {ιₛᵢ : Type} (OStmtIn : ιₛᵢ → Type) (WitIn : Type)
    (StmtOut : Type) {ιₛₒ : Type} (OStmtOut : ιₛₒ → Type) (WitOut : Type)
    {n : ℕ} (pSpec : ProtocolSpec n) :=
  Prover oSpec (StmtIn × (∀ i, OStmtIn i)) WitIn (StmtOut × (∀ i, OStmtOut i)) WitOut pSpec
\end{leancode}
\src{\file{ArkLib/}\allowbreak\file{OracleReduction/}\allowbreak\file{Basic.lean:255-260}, ArkLib \code{dca90385}}

\begin{define}{\lean{OracleVerifier}: a verifier that only queries}
\lean{verify} sees the statement and the challenges and may \emph{query} the input oracles and
the prover's messages (it never sees a message whole). \lean{outputOracle} says which oracles it
hands on: either an embedding naming, for each output oracle, an input oracle or a message it
\emph{is} (\lean{OracleOutputEmbedding}), or a simulation computing a virtual output oracle query
by query (\lean{OracleOutputSimulation}).
\end{define}
\begin{leancode}
@[ext]
structure OracleVerifier {ι : Type} (oSpec : OracleSpec ι)
    (StmtIn : Type) {ιₛᵢ : Type} (OStmtIn : ιₛᵢ → Type)
    (StmtOut : Type) {ιₛₒ : Type} (OStmtOut : ιₛₒ → Type)
    {n : ℕ} (pSpec : ProtocolSpec n)
    [Oₛᵢ : ∀ i, OracleInterface (OStmtIn i)]
    [Oₘ : ∀ i, OracleInterface (pSpec.Message i)]
    [Oₛₒ : ∀ i, OracleInterface (OStmtOut i)]
    where

  /-- The core verification logic. Takes the input statement `stmtIn` and all verifier challenges
  `challenges` (which are determined outside this function, typically by sampling for
  public-coin protocols). Returns the output statement `StmtOut` within an `OracleComp` that has
  access to external oracles `oSpec`, input statement oracles `OStmtIn`, and prover message
  oracles `pSpec.Message`. -/
  verify : StmtIn → pSpec.Challenges →
    OptionT (OracleComp (oSpec + ([OStmtIn]ₒ + [pSpec.Message]ₒ))) StmtOut

  /-- The output family has exactly one semantic representation: either the
  legacy embedded-source form or a virtual query implementation. -/
  outputOracle :
    OracleOutputEmbedding OStmtIn pSpec.Message OStmtOut ⊕
      OracleOutputSimulation oSpec OStmtIn OStmtOut pSpec
\end{leancode}
\src{\file{ArkLib/}\allowbreak\file{OracleReduction/}\allowbreak\file{Basic.lean:322-344}, ArkLib \code{dca90385}}

\begin{leancode}
structure OracleOutputEmbedding {ιₛᵢ ιₛₒ : Type}
    (OStmtIn : ιₛᵢ → Type) {ιₘ : Type} (Message : ιₘ → Type)
    (OStmtOut : ιₛₒ → Type)
    [Oₛᵢ : ∀ i, OracleInterface (OStmtIn i)]
    [Oₘ : ∀ i, OracleInterface (Message i)]
    [Oₛₒ : ∀ i, OracleInterface (OStmtOut i)] where
  embed : ιₛₒ ↪ ιₛᵢ ⊕ ιₘ
  hEq : ∀ i, OStmtOut i = match embed i with
    | Sum.inl j => OStmtIn j
    | Sum.inr j => Message j
  outputInterface_heq : ∀ i, match embed i with
    | Sum.inl j => HEq (Oₛₒ i) (Oₛᵢ j)
    | Sum.inr j => HEq (Oₛₒ i) (Oₘ j)
\end{leancode}
\src{\file{ArkLib/}\allowbreak\file{OracleReduction/}\allowbreak\file{Basic.lean:295-307}, ArkLib \code{dca90385}}

\begin{define}{\lean{toVerifier}: the plain verifier of an oracle verifier}
It answers the oracle verifier's queries from the actual oracle contents and messages of the
transcript, and materialises the output oracles. Every security notion of an oracle verifier is
the notion of its \lean{toVerifier} (\cref{sec:gen-lib-ark-toverifier}).
\end{define}
\begin{leancode}
def toVerifier : Verifier oSpec (StmtIn × ∀ i, OStmtIn i) (StmtOut × (∀ i, OStmtOut i)) pSpec where
  verify := fun ⟨stmt, oStmt⟩ transcript => OptionT.mk <|
    Option.map (fun stmtOut =>
      (stmtOut, verifier.materializeOutput
        transcript.challenges oStmt transcript.messages)) <$>
      simulateQ (OracleInterface.simOracle2 oSpec oStmt transcript.messages)
        (verifier.verify stmt transcript.challenges).run
\end{leancode}
\src{\file{ArkLib/}\allowbreak\file{OracleReduction/}\allowbreak\file{Basic.lean:490-496}, ArkLib \code{dca90385}}

\begin{define}{\lean{Reduction}, \lean{OracleReduction}, \lean{OracleProof}}
A reduction is a prover and a verifier; an oracle reduction is an oracle prover and an oracle
verifier; an oracle proof is an oracle reduction whose output is a Boolean, with no output oracle
and no output witness.
\end{define}
\begin{leancode}
@[ext]
structure Reduction {ι : Type} (oSpec : OracleSpec ι)
    (StmtIn WitIn StmtOut WitOut : Type) {n : ℕ} (pSpec : ProtocolSpec n) where
  prover : Prover oSpec StmtIn WitIn StmtOut WitOut pSpec
  verifier : Verifier oSpec StmtIn StmtOut pSpec
\end{leancode}
\src{\file{ArkLib/}\allowbreak\file{OracleReduction/}\allowbreak\file{Basic.lean:624-628}, ArkLib \code{dca90385}}

\begin{leancode}
@[ext]
structure OracleReduction {ι : Type} (oSpec : OracleSpec ι)
    (StmtIn : Type) {ιₛᵢ : Type} (OStmtIn : ιₛᵢ → Type) (WitIn : Type)
    (StmtOut : Type) {ιₛₒ : Type} (OStmtOut : ιₛₒ → Type) (WitOut : Type)
    {n : ℕ} (pSpec : ProtocolSpec n)
    [Oₛᵢ : ∀ i, OracleInterface (OStmtIn i)] [Oₘ : ∀ i, OracleInterface (pSpec.Message i)]
    [Oₛₒ : ∀ i, OracleInterface (OStmtOut i)]
    where
  prover : OracleProver oSpec StmtIn OStmtIn WitIn StmtOut OStmtOut WitOut pSpec
  verifier : OracleVerifier oSpec StmtIn OStmtIn StmtOut OStmtOut pSpec
\end{leancode}
\src{\file{ArkLib/}\allowbreak\file{OracleReduction/}\allowbreak\file{Basic.lean:632-641}, ArkLib \code{dca90385}}

\begin{leancode}
@[reducible] def OracleProof {ι : Type} (oSpec : OracleSpec ι)
    (Statement : Type) {ιₛᵢ : Type} (OStatement : ιₛᵢ → Type) (Witness : Type)
    {n : ℕ} (pSpec : ProtocolSpec n)
    [Oₛᵢ : ∀ i, OracleInterface (OStatement i)]
    [Oₘ : ∀ i, OracleInterface (pSpec.Message i)] :=
  @OracleReduction ι oSpec Statement ιₛᵢ OStatement Witness Bool Empty
    (fun _ : Empty => Unit) Unit n pSpec Oₛᵢ Oₘ (fun i => nomatch i)
\end{leancode}
\src{\file{ArkLib/}\allowbreak\file{OracleReduction/}\allowbreak\file{Basic.lean:669-675}, ArkLib \code{dca90385}}

\subsection{Running a protocol}\label{sec:gen-lib-ark-run}

\begin{define}{A run}
\lean{Prover.runToRound} folds the rounds: at a challenge round the challenge is obtained by a
query to the challenge oracle \lean{[pSpec.Challenge]ₒ} (one oracle per challenge round, trivial
query) and handed to \lean{receiveChallenge}; at a message round \lean{sendMessage} runs.
\lean{Reduction.run} runs the prover, then the verifier on the transcript; the verifier's failure
makes the run fail.
\end{define}
\begin{leancode}
def processRound (j : Fin n)
    (prover : Prover oSpec StmtIn WitIn StmtOut WitOut pSpec)
    (currentResult : OracleComp (oSpec + [pSpec.Challenge]ₒ)
      (pSpec.Transcript j.castSucc × prover.PrvState j.castSucc)) :
      OracleComp (oSpec + [pSpec.Challenge]ₒ)
        (pSpec.Transcript j.succ × prover.PrvState j.succ) := do
  let ⟨transcript, state⟩ ← currentResult
  match hDir : pSpec.dir j with
  | .V_to_P => do
    let challenge ← pSpec.getChallenge ⟨j, hDir⟩
    letI newState := (← prover.receiveChallenge ⟨j, hDir⟩ state) challenge
    return ⟨transcript.concat challenge, newState⟩
  | .P_to_V => do
    let ⟨msg, newState⟩ ← prover.sendMessage ⟨j, hDir⟩ state
    return ⟨transcript.concat msg, newState⟩
\end{leancode}
\src{\file{ArkLib/}\allowbreak\file{OracleReduction/}\allowbreak\file{Execution.lean:96-110}, ArkLib \code{dca90385}}

\begin{leancode}
def Reduction.run (stmt : StmtIn) (wit : WitIn)
    (reduction : Reduction oSpec StmtIn WitIn StmtOut WitOut pSpec) :
      OptionT (OracleComp (oSpec + [pSpec.Challenge]ₒ))
        ((FullTranscript pSpec × StmtOut × WitOut) × StmtOut) := do
  -- `ctxOut` contains both the output statement and witness after running the prover
  let proverResult ← reduction.prover.run stmt wit
  let stmtOut ← liftM (reduction.verifier.run stmt proverResult.1).run
  return ⟨proverResult, ← stmtOut.getM⟩
\end{leancode}
\src{\file{ArkLib/}\allowbreak\file{OracleReduction/}\allowbreak\file{Execution.lean:209-216}, ArkLib \code{dca90385}}

In the security definitions the challenge oracle is implemented by uniform sampling
(\lean{challengeQueryImpl}) and the shared oracle \lean{oSpec} by
\lean{impl : QueryImpl oSpec (StateT σ ProbComp)} from an initial state drawn by
\lean{init : ProbComp σ}. leanerVM's shared oracle is empty (\lean{[]ₒ}), so \lean{init} and
\lean{impl} are inert and the master theorems quantify over them.
\begin{leancode}
def challengeQueryImpl {pSpec : ProtocolSpec n} [∀ i, SampleableType (pSpec.Challenge i)] :
    QueryImpl ([pSpec.Challenge]ₒ'challengeOracleInterface) ProbComp :=
  fun q => $ᵗ (pSpec.Challenge q.1)
\end{leancode}
\src{\file{ArkLib/}\allowbreak\file{OracleReduction/}\allowbreak\file{ProtocolSpec/}\allowbreak\file{Basic.lean:732-734}, ArkLib \code{dca90385}}

\subsection{Security notions}\label{sec:gen-lib-ark-notions}

Completeness is in \cref{sec:gen-lib-ark-completeness}; round-by-round knowledge soundness, the
notion of the master theorem, is in \cref{sec:gen-lib-ark-sec}. The others:

\begin{define}{\lean{soundness}}
For every prover and every statement outside the input language, the run ends in the output
language with probability at most the error; failure counts for the verifier.
\end{define}
\begin{leancode}
def soundness (langIn : Set StmtIn) (langOut : Set StmtOut)
    (verifier : Verifier oSpec StmtIn StmtOut pSpec)
    (soundnessError : ℝ≥0) : Prop :=
  ∀ WitIn WitOut : Type,
  ∀ witIn : WitIn,
  ∀ prover : Prover oSpec StmtIn WitIn StmtOut WitOut pSpec,
  ∀ stmtIn ∉ langIn,
    let pImpl : QueryImpl (oSpec + [pSpec.Challenge]ₒ) (StateT σ ProbComp) :=
      impl.addLift challengeQueryImpl
    letI reduction := Reduction.mk prover verifier
    Pr[fun ⟨_, stmtOut⟩ => stmtOut ∈ langOut | OptionT.mk do
      (simulateQ pImpl (reduction.run stmtIn witIn).run).run' (← init)] ≤ soundnessError
\end{leancode}
\src{\file{ArkLib/}\allowbreak\file{OracleReduction/}\allowbreak\file{Security/}\allowbreak\file{Basic.lean:272-283}, ArkLib \code{dca90385}}

\begin{define}{\lean{Extractor.Straightline}, \lean{knowledgeSoundness}}
A straight-line extractor gets the statement, the output witness, the transcript and both query
logs and returns an input witness (or fails). Knowledge soundness bounds the probability that the
run's output is in the output relation while the extracted input witness is not in the input
relation (an extractor that fails counts against it, \file{:329-342}).
\end{define}
\begin{leancode}
def Straightline :=
  StmtIn → -- input statement
  WitOut → -- output witness
  FullTranscript pSpec → -- reduction transcript
  QueryLog oSpec → -- prover's query log
  QueryLog oSpec → -- verifier's query log
  OptionT (OracleComp oSpec) WitIn -- input witness
\end{leancode}
\src{\file{ArkLib/}\allowbreak\file{OracleReduction/}\allowbreak\file{Security/}\allowbreak\file{Basic.lean:248-254}, ArkLib \code{dca90385}}

\begin{leancode}
def knowledgeSoundness (relIn : Set (StmtIn × WitIn)) (relOut : Set (StmtOut × WitOut))
    (verifier : Verifier oSpec StmtIn StmtOut pSpec) (knowledgeError : ℝ≥0) : Prop :=
  ∃ extractor : Extractor.Straightline oSpec StmtIn WitIn WitOut pSpec,
  ∀ stmtIn : StmtIn,
  ∀ witIn : WitIn,
  ∀ prover : Prover oSpec StmtIn WitIn StmtOut WitOut pSpec,
    let pImpl : QueryImpl (oSpec + [pSpec.Challenge]ₒ) (StateT σ ProbComp) :=
      impl.addLift challengeQueryImpl
    let exec := do
      let ⟨⟨⟨transcript, ⟨_, witOut⟩⟩, stmtOut⟩, proveQueryLog, verifyQueryLog⟩
        ← (Reduction.mk prover verifier).runWithLog stmtIn witIn
      let extractedWitIn? ←
        liftM (extractor stmtIn witOut transcript proveQueryLog.fst verifyQueryLog).run
      return (stmtIn, extractedWitIn?, stmtOut, witOut)
    Pr[fun ⟨stmtIn, extractedWitIn?, stmtOut, witOut⟩ =>
        (∀ extractedWitIn ∈ extractedWitIn?, (stmtIn, extractedWitIn) ∉ relIn) ∧
          (stmtOut, witOut) ∈ relOut
      | OptionT.mk do (simulateQ pImpl exec.run).run' (← init)] ≤ knowledgeError
\end{leancode}
\src{\file{ArkLib/}\allowbreak\file{OracleReduction/}\allowbreak\file{Security/}\allowbreak\file{Basic.lean:344-361}, ArkLib \code{dca90385}}

\begin{define}{\lean{Verifier.StateFunction}, \lean{rbrSoundness}}
Round-by-round soundness (\cref{sec:gen-lib-ark-literature}): a predicate on partial transcripts,
true at the start exactly on the language, preserved by prover messages when false, false at the
end unless the verifier can output a statement of the output language, and made true by a
challenge with small probability. \lean{rbrSoundnessWorstCase}
(\file{Security/}\allowbreak\file{RoundByRound.lean:517-527}) is its per-prefix form.
\end{define}
\begin{leancode}
structure StateFunction
    (langIn : Set StmtIn) (langOut : Set StmtOut)
    (verifier : Verifier oSpec StmtIn StmtOut pSpec)
    where
  toFun : (m : Fin (n + 1)) → StmtIn → Transcript m pSpec → Prop
  /-- For all input statement not in the language, the state function is false for that statement
    and the empty transcript -/
  toFun_empty : ∀ stmt, stmt ∈ langIn ↔ toFun 0 stmt default
  /-- If the state function is false for a partial transcript, and the next message is from the
    prover to the verifier, then the state function is also false for the new partial transcript
    regardless of the message -/
  toFun_next : ∀ m, pSpec.dir m = .P_to_V →
    ∀ stmt tr, ¬ toFun m.castSucc stmt tr →
    ∀ msg, ¬ toFun m.succ stmt (tr.concat msg)
  /-- If the state function is false for a full transcript, the verifier will not output a statement
    in the output language -/
  toFun_full : ∀ stmt tr, ¬ toFun (.last n) stmt tr →
    Pr[(· ∈ langOut) | OptionT.mk do (simulateQ impl (verifier.run stmt tr)).run' (← init)] = 0
\end{leancode}
\src{\file{ArkLib/}\allowbreak\file{OracleReduction/}\allowbreak\file{Security/}\allowbreak\file{RoundByRound.lean:135-152}, ArkLib \code{dca90385}}

\begin{leancode}
def rbrSoundness (langIn : Set StmtIn) (langOut : Set StmtOut)
    (verifier : Verifier oSpec StmtIn StmtOut pSpec)
    (rbrSoundnessError : pSpec.ChallengeIdx → ℝ≥0) : Prop :=
  ∃ stateFunction : verifier.StateFunction init impl langIn langOut,
  ∀ stmtIn ∉ langIn,
  ∀ WitIn WitOut : Type,
  ∀ witIn : WitIn,
  ∀ prover : Prover oSpec StmtIn WitIn StmtOut WitOut pSpec,
  ∀ i : pSpec.ChallengeIdx,
    Pr[fun ⟨transcript, challenge⟩ =>
      ¬ stateFunction i.1.castSucc stmtIn transcript ∧
        stateFunction i.1.succ stmtIn (transcript.concat challenge)
    | do
      (simulateQ (impl.addLift challengeQueryImpl : QueryImpl _ (StateT σ ProbComp))
        (do
          let ⟨transcript, _⟩ ← prover.runToRound i.1.castSucc stmtIn witIn
          let challenge ← liftComp (pSpec.getChallenge i) _
          return (transcript, challenge))).run' (← init)] ≤
      rbrSoundnessError i
\end{leancode}
\src{\file{ArkLib/}\allowbreak\file{OracleReduction/}\allowbreak\file{Security/}\allowbreak\file{RoundByRound.lean:347-365}, ArkLib \code{dca90385}}

\begin{define}{\lean{KnowledgeStateFunction}, \lean{Extractor.RoundByRound} and the round-by-round knowledge notions}
Defined and discussed in \cref{sec:gen-lib-ark-definition}. The existential worst-case form
quantifies the intermediate witness types, the extractor and the knowledge state function
existentially:
\end{define}
\begin{leancode}
def rbrKnowledgeSoundnessWorstCase (relIn : Set (StmtIn × WitIn))
    (relOut : Set (StmtOut × WitOut))
    (verifier : Verifier oSpec StmtIn StmtOut pSpec)
    (rbrKnowledgeError : pSpec.ChallengeIdx → ℝ≥0) : Prop :=
  ∃ WitMid : Fin (n + 1) → Type,
  ∃ extractor : Extractor.RoundByRound oSpec StmtIn WitIn WitOut pSpec WitMid,
  ∃ kSF : verifier.KnowledgeStateFunction init impl relIn relOut extractor,
  ∀ stmtIn : StmtIn,
  ∀ i : pSpec.ChallengeIdx,
  ∀ transcript : Transcript i.1.castSucc pSpec,
    Pr[fun challenge =>
      ∃ witMid,
        ¬ kSF i.1.castSucc stmtIn transcript
          (extractor.extractMid i.1 stmtIn (transcript.concat challenge) witMid) ∧
          kSF i.1.succ stmtIn (transcript.concat challenge) witMid
      | $ᵗ (pSpec.Challenge i)] ≤ rbrKnowledgeError i
\end{leancode}
\src{\file{ArkLib/}\allowbreak\file{OracleReduction/}\allowbreak\file{Security/}\allowbreak\file{RoundByRound.lean:534-549}, ArkLib \code{dca90385}}

\begin{define}{\lean{Verifier.GuardedForm}, \lean{Prover.OutputIsPure}}
A guarded verifier equals \enquote{if \lean{check} then \lean{pure (out …)} else fail}, with
\lean{check} and \lean{out} as data. A prover's output is pure when it is a function of its final
state with no oracle query.
\end{define}
\begin{leancode}
def IsGuardedWith (V : Verifier oSpec StmtIn StmtOut pSpec)
    (check : StmtIn → FullTranscript pSpec → Bool)
    (out : StmtIn → FullTranscript pSpec → StmtOut) : Prop :=
  ∀ stmt tr, V.verify stmt tr = if check stmt tr then pure (out stmt tr) else failure
\end{leancode}
\src{\file{ArkLib/}\allowbreak\file{OracleReduction/}\allowbreak\file{Security/}\allowbreak\file{CoordinateWiseSpecialSoundness/}\allowbreak\file{Guarded.lean:86-89}, ArkLib \code{dca90385}}

\begin{leancode}
structure GuardedForm (V : Verifier oSpec StmtIn StmtOut pSpec) where
  /-- The runtime guard. -/
  check : StmtIn → FullTranscript pSpec → Bool
  /-- The verdict where the guard passes. -/
  out : StmtIn → FullTranscript pSpec → StmtOut
  /-- The verifier is guarded with exactly these. -/
  verify_eq : V.IsGuardedWith check out
\end{leancode}
\src{\file{ArkLib/}\allowbreak\file{OracleReduction/}\allowbreak\file{Security/}\allowbreak\file{CoordinateWiseSpecialSoundness/}\allowbreak\file{Guarded.lean:112-118}, ArkLib \code{dca90385}}

\begin{leancode}
/-- The prover's output is a deterministic function of its final private state, with no
oracle queries. This condition does not constrain the message-sending steps. -/
class Prover.OutputIsPure (P : Prover oSpec StmtIn WitIn StmtOut WitOut pSpec) where
    output_is_pure : ∃ output : _ → _, ∀ st, P.output st = pure (output st)
\end{leancode}
\src{\file{ArkLib/}\allowbreak\file{OracleReduction/}\allowbreak\file{Basic.lean:989-992}, ArkLib \code{dca90385}}

\subsection{Composition}\label{sec:gen-lib-ark-append}

\begin{define}{\lean{append} of provers, verifiers and reductions}
\lean{Prover.append} (\file{Append/}\allowbreak\file{Basic.lean:58-148}) runs the first prover and feeds its output
to the second prover's \lean{input} at the seam; \lean{Verifier.append} runs the first verifier on
the first half of the transcript and the second on the second half at the first's output;
\lean{OracleVerifier.append} (\file{:619-629}) routes the second verifier's queries through the
first's output-oracle simulation; \lean{OracleReduction.append} pairs them; \lean{seqCompose}
(\file{General.lean:255-268}) iterates.
\end{define}
\begin{leancode}
def Verifier.append (V₁ : Verifier oSpec Stmt₁ Stmt₂ pSpec₁)
    (V₂ : Verifier oSpec Stmt₂ Stmt₃ pSpec₂) :
      Verifier oSpec Stmt₁ Stmt₃ (pSpec₁ ++ₚ pSpec₂) where
  verify := fun stmt transcript => do
    return ← V₂.verify (← V₁.verify stmt transcript.fst) transcript.snd

/-- Compose the component provers and verifiers of two reductions. -/
def Reduction.append (R₁ : Reduction oSpec Stmt₁ Wit₁ Stmt₂ Wit₂ pSpec₁)
    (R₂ : Reduction oSpec Stmt₂ Wit₂ Stmt₃ Wit₃ pSpec₂) :
      Reduction oSpec Stmt₁ Wit₁ Stmt₃ Wit₃ (pSpec₁ ++ₚ pSpec₂) where
  prover := Prover.append R₁.prover R₂.prover
  verifier := Verifier.append R₁.verifier R₂.verifier
\end{leancode}
\src{\file{ArkLib/}\allowbreak\file{OracleReduction/}\allowbreak\file{Composition/}\allowbreak\file{Sequential/}\allowbreak\file{Append/}\allowbreak\file{Basic.lean:236-247}, ArkLib \code{dca90385}}

\begin{leancode}
def OracleReduction.append (R₁ : OracleReduction oSpec Stmt₁ OStmt₁ Wit₁ Stmt₂ OStmt₂ Wit₂ pSpec₁)
    (R₂ : OracleReduction oSpec Stmt₂ OStmt₂ Wit₂ Stmt₃ OStmt₃ Wit₃ pSpec₂) :
      OracleReduction oSpec Stmt₁ OStmt₁ Wit₁ Stmt₃ OStmt₃ Wit₃ (pSpec₁ ++ₚ pSpec₂) where
  prover := Prover.append R₁.prover R₂.prover
  verifier := OracleVerifier.append R₁.verifier R₂.verifier
\end{leancode}
\src{\file{ArkLib/}\allowbreak\file{OracleReduction/}\allowbreak\file{Composition/}\allowbreak\file{Sequential/}\allowbreak\file{Append/}\allowbreak\file{Basic.lean:709-713}, ArkLib \code{dca90385}}

\begin{define}{\lean{Extractor.}\allowbreak\lean{RoundByRound.}\allowbreak\lean{append}, \lean{Verifier.}\allowbreak\lean{StateFunction.}\allowbreak\lean{append}}
The appended extractor runs the second extractor on the second half at the statement
\lean{verify s tr₁} that a \emph{given} verdict function assigns, then the first on the first
half; at the seam it hands the second extractor's round-0 witness to the first's
\lean{extractOut} (\file{Append/}\allowbreak\file{StateFunction.lean:75-139}). The appended (language) state
function is disjunctive past the seam (\file{:292}, docstring \file{:17-19}). The knowledge
version is leanerVM's port (\cref{sec:gen-lib-ark-port}).
\end{define}
\begin{leancode}
/-- Compose round-by-round extractors using a deterministic intermediate-statement function.
The composed witness family retains the first extractor's final witness at the boundary. -/
def RoundByRound.append
    {WitMid₁ : Fin (m + 1) → Type} {WitMid₂ : Fin (n + 1) → Type}
    (E₁ : Extractor.RoundByRound oSpec Stmt₁ Wit₁ Wit₂ pSpec₁ WitMid₁)
    (E₂ : Extractor.RoundByRound oSpec Stmt₂ Wit₂ Wit₃ pSpec₂ WitMid₂)
    (verify : Stmt₁ → pSpec₁.FullTranscript → Stmt₂) :
      Extractor.RoundByRound oSpec Stmt₁ Wit₁ Wit₃ (pSpec₁ ++ₚ pSpec₂)
        (Fin.append (m := m + 1) WitMid₁ (Fin.tail WitMid₂) ∘ Fin.cast (by omega)) where
\end{leancode}
\src{\file{ArkLib/}\allowbreak\file{OracleReduction/}\allowbreak\file{Composition/}\allowbreak\file{Sequential/}\allowbreak\file{Append/}\allowbreak\file{StateFunction.lean:73-81}, ArkLib \code{dca90385}}

The composition of completeness (\lean{append_}\allowbreak\lean{perfectCompleteness_}\allowbreak\lean{of_guarded_}\allowbreak\lean{verifiers},
\lean{append_}\allowbreak\lean{completeness_}\allowbreak\lean{of_guarded_}\allowbreak\lean{verifiers}, \lean{seqCompose_*}) is in
\cref{sec:gen-lib-ark-completeness}; all are proved at both pins.

\subsection{Fiat--Shamir and commitments}\label{sec:gen-lib-ark-fs}

\begin{define}{\lean{Reduction.fiatShamir}, \lean{Verifier.fiatShamir}, the challenge oracle}
The \enquote{slow} Fiat--Shamir transformation: the challenge of round $i$ is the answer of an
oracle whose query is the statement together with all messages before round $i$
(\lean{challengeOracleInterfaceSR}); the transformed verifier receives all messages in one round,
derives the challenges by querying that oracle and runs the original verifier.
\end{define}
\begin{leancode}
@[reducible, inline, specialize]
def challengeOracleInterfaceSR (StmtIn : Type) (pSpec : ProtocolSpec n) :
    ∀ i, OracleInterface (pSpec.Challenge i) := fun i =>
  { Query := StmtIn × pSpec.MessagesUpTo i.1.castSucc
    toOC.spec := fun _ => pSpec.Challenge i
    toOC.impl := fun _ => read }
\end{leancode}
\src{\file{ArkLib/}\allowbreak\file{OracleReduction/}\allowbreak\file{ProtocolSpec/}\allowbreak\file{Basic.lean:829-834}, ArkLib \code{dca90385}}

\begin{leancode}
/-- The (slow) Fiat-Shamir transformation for the verifier. -/
def Verifier.fiatShamir (V : Verifier oSpec StmtIn StmtOut pSpec) :
    NonInteractiveVerifier (∀ i, pSpec.Message i) (oSpec + fsChallengeOracle StmtIn pSpec)
      StmtIn StmtOut where
  verify := fun stmtIn proof => do
    let messages : pSpec.Messages := proof 0
    let transcript ← (messages.deriveTranscriptFS (oSpec := oSpec) stmtIn)
    Option.getM (← (V.verify stmtIn transcript).run)

/-- The Fiat-Shamir transformation for an (interactive) reduction, which consists of applying the
  Fiat-Shamir transformation to both the prover and the verifier. -/
def Reduction.fiatShamir (R : Reduction oSpec StmtIn WitIn StmtOut WitOut pSpec) :
    NonInteractiveReduction (∀ i, pSpec.Message i) (oSpec + fsChallengeOracle StmtIn pSpec)
      StmtIn WitIn StmtOut WitOut where
  prover := R.prover.fiatShamir
  verifier := R.verifier.fiatShamir
\end{leancode}
\src{\file{ArkLib/}\allowbreak\file{OracleReduction/}\allowbreak\file{FiatShamir/}\allowbreak\file{Basic.lean:129-144}, ArkLib \code{dca90385}}

\begin{define}{\lean{Commitment.Scheme}, \lean{binding}, \lean{extractability}}
Keys, \lean{commit : ComKey → Data → OracleComp (Commitment × Decommitment)}, and an opening
\emph{proof} for the statement \enquote{commitment, query, response} with witness \enquote{data,
decommitment}. \lean{binding} bounds any adversary's probability of opening one commitment to two
different responses at the same query; \lean{extractability} is a placeholder (its body is
\lean{False}, at both pins: \file{Commitments/}\allowbreak\file{Functional/}\allowbreak\file{Basic.lean:250-255} at \code{dca90385},
\file{:249-254} at \code{7653a901}).
\end{define}
\begin{leancode}
/-- Key generation for a commitment scheme, producing a committer key and a verifier key. -/
structure KeyGen where
  keygen : OracleComp oSpec (ComKey × VerifKey)

/-- The commitment algorithm, parameterized by the committer key and the data to commit. -/
structure Commit where
  commit : ComKey → Data → OracleComp oSpec (Commitment × Decommitment)

variable [O : OracleInterface Data] {n : ℕ} (pSpec : ProtocolSpec n)

/-- The opening protocol used to prove a claimed oracle response for committed data. -/
structure Opening where
  opening : (ComKey × VerifKey) →
    Proof oSpec (Commitment × (q : O.Query) × O.Response q) (Data × Decommitment) pSpec

/-- A commitment scheme with key generation, commitment, and opening algorithms. -/
structure Scheme extends
    KeyGen oSpec ComKey VerifKey,
    Commit oSpec Data Commitment Decommitment ComKey,
    Opening oSpec Data Commitment Decommitment ComKey VerifKey pSpec
\end{leancode}
\src{\file{ArkLib/}\allowbreak\file{Commitments/}\allowbreak\file{Functional/}\allowbreak\file{Basic.lean:51-70}, ArkLib \code{dca90385}}

\begin{leancode}
def binding (bindingError : ℝ≥0) : Prop :=
  ∀ AuxState : Type,
  ∀ adversary : BindingAdversary oSpec Data Commitment AuxState pSpec ComKey,
    bindingExperiment init impl scheme AuxState adversary ≤ bindingError
\end{leancode}
\src{\file{ArkLib/}\allowbreak\file{Commitments/}\allowbreak\file{Functional/}\allowbreak\file{Basic.lean:220-223}, ArkLib \code{dca90385}}

\subsection{Generic components and the sumcheck specification}\label{sec:gen-lib-ark-components}

\begin{define}{Generic components}
\lean{DoNothing}: the identity reduction (\file{ProofSystem/}\allowbreak\file{Component/}\allowbreak\file{DoNothing.lean:33-66}).
\lean{ReduceClaim}: zero rounds, maps the statement and the witness (\file{ReduceClaim.lean:50-66};
oracle version \file{:282-323}). \lean{CheckClaim}: zero rounds; the plain verifier
\lean{guard}s a decidable predicate (\file{CheckClaim.lean:72-73}); the oracle verifier is a
pass-through (\file{:228-231}). \lean{RandomQuery}: one challenge, a query; the verifier outputs
the query and both oracles; the output relation says the two oracles agree there
(\file{RandomQuery.lean:45-57}, \file{:94-111}). \lean{SendSingleWitness}: one prover message, the
witness, which becomes an oracle (\file{SendWitness.lean:316-343}; completeness admitted,
\file{:440-443}).
\end{define}
\begin{leancode}
/-- The verifier for the `ReduceClaim` reduction. -/
def verifier : Verifier oSpec StmtIn StmtOut !p[] where
  verify := fun stmt _ => pure (mapStmt stmt)

/-- The reduction for the `ReduceClaim` reduction. -/
def reduction : Reduction oSpec StmtIn WitIn StmtOut WitOut !p[] where
\end{leancode}
\src{\file{ArkLib/}\allowbreak\file{ProofSystem/}\allowbreak\file{Component/}\allowbreak\file{ReduceClaim.lean:61-66}, ArkLib \code{dca90385}}

\begin{leancode}
@[inline, specialize]
def verifier : Verifier oSpec Statement Statement !p[] where
  verify := fun stmt _ => do guard (pred stmt); return stmt
\end{leancode}
\src{\file{ArkLib/}\allowbreak\file{ProofSystem/}\allowbreak\file{Component/}\allowbreak\file{CheckClaim.lean:71-73}, ArkLib \code{dca90385}}

\begin{leancode}
@[reducible, simp]
def relIn : Set ((StmtIn × ∀ i, OStmtIn OStatement i) × WitIn) :=
  { ⟨⟨(), oracles⟩, ()⟩ | oracles 0 = oracles 1 }

/--
The output relation states that if the verifier's single query was `q`, then
`a` and `b` agree on that `q`, i.e. `answer a q = answer b q`.
-/
@[reducible, simp]
def relOut : Set ((StmtOut OStatement × ∀ i, OStmtOut OStatement i) × WitOut) :=
  { ⟨⟨q, oStmt⟩, ()⟩ | answer (oStmt 0) q = answer (oStmt 1) q }

@[reducible]
def pSpec : ProtocolSpec 1 := ⟨!v[.V_to_P], !v[Query OStatement]⟩
\end{leancode}
\src{\file{ArkLib/}\allowbreak\file{ProofSystem/}\allowbreak\file{Component/}\allowbreak\file{RandomQuery.lean:44-57}, ArkLib \code{dca90385}}

\begin{define}{The sumcheck specification}
A round's statement is the running target and the challenges so far; the oracle statement is a
multivariate polynomial of individual degree at most \lean{deg} in $n$ variables; the round
relation says the sum of the polynomial over the remaining cube ($D^{n-i}$) at the challenges
equals the target; a round is \enquote{prover sends a univariate of degree at most \lean{deg},
verifier sends a field element}; the whole protocol is the \lean{seqCompose} of $n$ rounds. Both
the completeness and the round-by-round knowledge soundness of the specification are admitted at
both pins (\cref{sec:gen-lib-ark-baseline}). \lean{SumcheckDomain} (\file{Domain.lean:55-59})
generalises the evaluation domain to a per-coordinate family.
\end{define}
\begin{leancode}
structure StatementRound (i : Fin (n + 1)) where
  -- The target value for sum-check
  target : R
  -- The challenges sent from the verifier to the prover from previous rounds
  challenges : Fin i → R

abbrev OutputStatement := StatementRound R _ (.last n)

/-- Oracle statement for sum-check, which is a multivariate polynomial over `n` variables of
  individual degree at most `deg`, equipped with the poly evaluation oracle interface. -/
@[reducible]
def OracleStatement : Unit → Type := fun _ => R⦃≤ deg⦄[X Fin n]

/-- The sum-check relation for the `i`-th round, for `i ≤ n` -/
def relationRound (i : Fin (n + 1)) :
    Set (((StatementRound R n i) × (∀ i, OracleStatement R n deg i)) × Unit) :=
\end{leancode}
\src{\file{ArkLib/}\allowbreak\file{ProofSystem/}\allowbreak\file{Sumcheck/}\allowbreak\file{Spec/}\allowbreak\file{SingleRound.lean:130-145}, ArkLib \code{dca90385}}

\begin{leancode}
namespace SingleRound

/-- The protocol specification for a single round of sum-check.
Has the form `⟨!v[.P_to_V, .V_to_P], !v[R⦃≤ deg⦄[X], R]⟩` -/
@[reducible]
def pSpec : ProtocolSpec 2 :=
\end{leancode}
\src{\file{ArkLib/}\allowbreak\file{ProofSystem/}\allowbreak\file{Sumcheck/}\allowbreak\file{Spec/}\allowbreak\file{SingleRound.lean:149-154}, ArkLib \code{dca90385}}

\begin{leancode}
def pSpec : ProtocolSpec (Fin.vsum (fun _ : Fin n => 2)) :=
  ProtocolSpec.seqCompose (fun _ => SingleRound.pSpec R deg)
  -- n * 2
  -- fun i => if i % 2 = 0 then (.P_to_V, R⦃≤ d⦄[X]) else (.V_to_P, R)

\end{leancode}
\src{\file{ArkLib/}\allowbreak\file{ProofSystem/}\allowbreak\file{Sumcheck/}\allowbreak\file{Spec/}\allowbreak\file{General.lean:129-133}, ArkLib \code{dca90385}}

\begin{leancode}

/-- Perfect completeness for the (full) sum-check protocol -/
theorem reduction_perfectCompleteness :
    (reduction R deg D n oSpec).perfectCompleteness init impl
      (relationRound R n deg D 0) (relationRound R n deg D (.last n)) :=
  Reduction.seqCompose_perfectCompleteness_of_guarded_verifiers
    (fun i => StatementRound R n i × ∀ j, OracleStatement R n deg j)
    (fun _ => Unit) init impl (relationRound R n deg D)
    (SingleRound.reduction R n deg D oSpec)
    (fun _ => inferInstance) (SingleRound.verifierGuardedForm R n deg D oSpec)
    (fun i s => SingleRound.reduction_perfectCompleteness (init := pure s) i)

/-- Round-by-round knowledge soundness with error `deg / |R|` per challenge for the (full)
  sum-check protocol -/
theorem oracleVerifier_rbrKnowledgeSoundness [Fintype R] :
    (oracleVerifier R deg D n oSpec).rbrKnowledgeSoundness init impl
      (relationRound R n deg D 0) (relationRound R n deg D (.last n))
      (fun _ => (deg : ℝ≥0) / (Fintype.card R)) :=
  OracleVerifier.seqCompose_rbrKnowledgeSoundness
    (rel := relationRound R n deg D)
\end{leancode}
\src{\file{ArkLib/}\allowbreak\file{ProofSystem/}\allowbreak\file{Sumcheck/}\allowbreak\file{Spec/}\allowbreak\file{General.lean:209-228}, ArkLib \code{dca90385}}

\par\endgroup